\documentclass[11pt,a4paper]{article}
\usepackage[utf8]{inputenc}
\usepackage[T1]{fontenc}
\usepackage[margin=2.5cm]{geometry}
\usepackage{amsmath,amssymb}
\usepackage{siunitx}
\usepackage{upgreek}
\DeclareSIPrefix{\micro}{\ensuremath{\upmu}}{-6}
\usepackage{graphicx}
\usepackage{booktabs}
\usepackage{array}
\usepackage{tikz}
\usepackage{listings}
\usepackage{xcolor}
\usepackage{caption}
\usepackage[hidelinks]{hyperref}

\usepackage{fancyhdr}
\usepackage{lastpage}

\graphicspath{{../figures/}{images/}}

% DTU corporate colours (DTU-templates/external/dtucolours.tex)
\definecolor{dtured}{rgb}{0.6,0,0}
\definecolor{dtugrey}{rgb}{0.6,0.6,0.6}

% Running header and footer on every page after the front page
\pagestyle{fancy}
\fancyhf{}
\fancyhead[L]{\small 34840 Fundamentals of Acoustics and Noise Control}
\fancyhead[R]{\includegraphics[height=0.9cm]{dtu_logo.pdf}}
\fancyfoot[L]{\small Mads Rudolph, s246132}
\fancyfoot[R]{\small Page \thepage\ of \pageref*{LastPage}}
\renewcommand{\headrulewidth}{0.4pt}
\renewcommand{\headrule}{\hbox to\headwidth{\color{dtured}\leaders\hrule height \headrulewidth\hfill}}
\setlength{\headheight}{28pt}
\sisetup{per-mode=symbol, detect-all}

\lstset{
  language=Matlab,
  basicstyle=\ttfamily\footnotesize,
  commentstyle=\color{green!40!black},
  keywordstyle=\color{blue!70!black},
  stringstyle=\color{purple!70!black},
  breaklines=true,
  columns=fullflexible,
  frame=single,
  rulecolor=\color{gray!50},
  numbers=left,
  numberstyle=\tiny\color{gray},
  literate={µ}{{$\mu$}}1
}

\newcommand{\answer}[1]{\par\smallskip\noindent\fbox{\parbox{\dimexpr\linewidth-2\fboxsep-2\fboxrule}{#1}}\par\smallskip}

\begin{document}

%--------------------------------------------------------------------- front page
\begin{titlepage}
\begin{tikzpicture}[remember picture, overlay]
  \node[anchor=north east, inner sep=0] at ([xshift=-2.5cm, yshift=-2cm]current page.north east)
    {\includegraphics[height=3.2cm]{dtu_logo.pdf}};
  \fill[dtured] ([yshift=1.6cm]current page.south west) rectangle ([yshift=2.0cm]current page.south east);
\end{tikzpicture}

\vspace*{5.5cm}
\noindent{\color{dtured}\rule{\linewidth}{1.2pt}}\\[0.8cm]
{\huge\bfseries Individual Hand-in 2026}\\[0.4cm]
{\Large 34840 Fundamentals of Acoustics and Noise Control}\\[0.4cm]
{\color{dtured}\rule{\linewidth}{1.2pt}}

\vspace{2.5cm}
\noindent\begin{tabular}{@{}ll}
  \textbf{Author:}     & Mads Rudolph \\
  \textbf{Study no.:}  & s246132 \\[0.3cm]
  \textbf{Course:}     & 34840 Fundamentals of Acoustics and Noise Control, autumn 2026 \\
  \textbf{Deadline:}   & 9 October 2026 \\
\end{tabular}

\vfill
\noindent\textbf{DTU Electro}\\
Department of Electrical and Photonics Engineering\\
Technical University of Denmark
\vspace*{1.2cm}
\end{titlepage}
\setcounter{page}{1}

\noindent Constants used throughout, as given in the assignment: air
$\rho_0 = \SI{1.2}{kg/m^3}$, $c = \SI{343}{m/s}$; water
$\rho_w = \SI{1000}{kg/m^3}$, $c_w = \SI{1480}{m/s}$; reference pressure
$p_\text{ref} = \SI{20}{\micro Pa}$. All calculations are done in the MATLAB
scripts \texttt{problem1.m} and \texttt{problem2.m}, listed in the appendix.

%=====================================================================
\section*{Problem 1 --- Plane wave on a coral reef}
%=====================================================================

A plane wave travels in water and hits a large reef at normal incidence. The
incident pressure amplitude is $\hat p_i = \SI{3}{Pa}$ and the total
intensity in front of the reef is $I_t = \SI{2.1}{\micro W/m^2}$. I read
``amplitude'' as the peak value of the harmonic pressure.

\subsection*{Q1.1 --- Absorption coefficient}

The characteristic impedance of water is
\begin{equation}
  Z_0 = \rho_w c_w = 1000 \cdot 1480 = \SI{1.48e6}{Pa.s/m}.
\end{equation}
For a plane wave the time-averaged intensity follows from the peak amplitude
as $I = \hat p^2/(2\rho c)$, so the incident wave carries
\begin{equation}
  I_i = \frac{\hat p_i^2}{2\rho_w c_w} = \frac{3^2}{2\cdot\num{1.48e6}}
      = \SI{3.04}{\micro W/m^2}.
\end{equation}

In front of the reef the field is the sum of the incident wave and the
reflected wave $\hat p_r = R\,\hat p_i$. The two waves travel in opposite
directions, so their intensities subtract (the cross terms between them
carry no net power in a plane-wave field). The measured total intensity is
therefore the net power per unit area that goes \emph{into} the reef:
\begin{equation}
  I_t = I_i - I_r = I_i\left(1 - |R|^2\right).
\end{equation}
The absorption coefficient is by definition the fraction of the incident
power that is not reflected, $\alpha = 1 - |R|^2$, which gives
\begin{equation}
  \alpha = \frac{I_t}{I_i} = \frac{2.1}{3.04} = 0.691.
\end{equation}

\answer{$\alpha \approx 0.69$ --- the reef absorbs about 69\,\% of the incident sound power.}

\subsection*{Q1.2 --- Surface impedance}

The reflection factor at normal incidence is set by the surface impedance
$Z_s = p/u$ of the reef:
\begin{equation}
  R = \frac{Z_s - Z_0}{Z_s + Z_0}.
\end{equation}
From Q1.1 the magnitude is
\begin{equation}
  |R| = \sqrt{1-\alpha} = \sqrt{1 - 0.691} = 0.556.
\end{equation}
Because $Z_s$ is real, $R$ is real as well, so either $R = +0.556$ or
$R = -0.556$. Solving the expression for $R$ for $Z_s$:
\begin{equation}
  Z_s = Z_0\,\frac{1+R}{1-R}.
\end{equation}
\begin{align}
  R = +0.556:&\quad Z_s = \num{1.48e6}\cdot\frac{1.556}{0.444}
     = 3.51\,Z_0 = \SI{5.19e6}{Pa.s/m},\\
  R = -0.556:&\quad Z_s = \num{1.48e6}\cdot\frac{0.444}{1.556}
     = 0.285\,Z_0 = \SI{4.22e5}{Pa.s/m}.
\end{align}
Both values give exactly the same absorption coefficient, so the intensity
measurement alone cannot tell them apart. To decide, I use what the reef
physically is: a rigid structure of calcium carbonate, which is both denser
and stiffer than water. It is acoustically \emph{harder} than water,
$Z_s > Z_0$, which means a positive reflection factor (pressure maximum at
the surface). The soft solution, $Z_s < Z_0$, would require a surface that
is acoustically softer than the water itself, e.g.\ a gas-filled layer, and
is not expected for a reef. Measuring the phase of the reflected wave (or
whether the standing-wave pattern has a pressure maximum or minimum at the
surface) would settle it.

\answer{$Z_s \approx \SI{5.2e6}{Pa.s/m}$ ($\approx 3.5\,\rho_w c_w$, $R = +0.556$).
The mathematically possible soft solution is $\SI{4.2e5}{Pa.s/m}$ ($R = -0.556$).}

%=====================================================================
\section*{Problem 2 --- Music festival}
%=====================================================================

Each stage plays 6 one-hour shows a day between 14:00 and 02:00. I assume the
stages contribute nothing outside the shows (background noise from other
sources is not part of the festival's emission).

\subsection*{Q2.1 --- 24 h equivalent level for stage 2}

The equivalent continuous level is the energy (mean-square pressure)
average over the period $T$:
\begin{equation}
  L_{\text{eq},T} = 10\log_{10}\!\left(\frac{1}{T}\sum_i t_i\,10^{L_i/10}\right).
\end{equation}
Stage 2 gives $L_2 = \SI{60}{dB}$ for $t_\text{on} = \SI{6}{h}$ and nothing for
the remaining 18 hours, so
\begin{equation}
  L_{\text{eq},24\text{h}} = 10\log_{10}\!\left(\frac{6}{24}\,10^{60/10}\right)
  = 60 + 10\log_{10}\frac{6}{24} = 60 - 6.0 = \SI{54.0}{dB}.
\end{equation}
The result does not depend on when during the day the six shows take place.

\answer{$L_{\text{eq},24\text{h}} \approx \SI{54.0}{dB}$ re \SI{20}{\micro Pa} for stage 2.}

\subsection*{Q2.2 --- 24 h equivalent level for the whole festival}

The three stages play different, unrelated programmes, so their signals are
incoherent and the mean-square pressures add. Each stage contributes for 6
hours. Because energy is summed over the full 24 hours, it makes no
difference whether the shows on the different stages overlap in time or
not:
\begin{equation}
  L_{\text{eq},24\text{h}}
  = 10\log_{10}\!\left(\frac{6}{24}\left(10^{58/10} + 10^{60/10} + 10^{54/10}\right)\right).
\end{equation}
The individual contributions are $\SI{52.0}{dB}$ (stage 1), $\SI{54.0}{dB}$
(stage 2) and $\SI{48.0}{dB}$ (stage 3), and the total is

\answer{$L_{\text{eq},24\text{h}} \approx \SI{56.7}{dB}$ re \SI{20}{\micro Pa} for the entire festival.}

\noindent For comparison, when all three stages play at the same time the
level at the neighbour is $\SI{62.7}{dB}$; the 24-hour average is
$\SI{6}{dB}$ lower because sound is only emitted a quarter of the day.

\subsection*{Q2.3 --- Subwoofer interference at 20 Hz}

\paragraph{Geometry.} I place the origin at the centre of the stage front
with $x$ along the stage and $y$ towards the audience
(Figure~\ref{fig:geometry}). The two subwoofers stand at the stage ends,
$(\mp10, 0)$\,m, Front of House (FOH) is at $(0, 13)$\,m and the five
audience positions are at $x = -20, -10, 0, 10, 20$\,m on the line
$y = \SI{19}{m}$.

\begin{figure}[htbp]
\centering
\begin{tikzpicture}[scale=0.21]
  \fill[gray!25] (-10,-2.2) rectangle (10,-0.4);
  \node at (0,-1.3) {\footnotesize stage};
  \fill (-10,0) ++(-0.6,-0.6) rectangle ++(1.2,1.2);
  \fill (10,0) ++(-0.6,-0.6) rectangle ++(1.2,1.2);
  \node[below left] at (-10.5,0) {\footnotesize L $(-10,0)$};
  \node[below right] at (10.5,0) {\footnotesize R $(10,0)$};
  \draw[fill=white] (0,13) circle (0.6);
  \node[right] at (0.8,13) {\footnotesize FOH $(0,13)$};
  \foreach \x in {-20,-10,0,10,20} {
    \draw[fill=white] (\x,19) circle (0.6);
    \node[above] at (\x,19.6) {\footnotesize $\x$};
  }
  \draw[dashed,gray] (-10,0) -- (0,13) -- (10,0);
  \draw[dashed,blue!60] (-10,0) -- (-10,19);
  \draw[dashed,blue!60] (10,0) -- (-10,19);
  \node[blue!60!black,left] at (-10,9.5) {\footnotesize $r_L=19.0$};
  \node[blue!60!black] at (5.5,4.5) {\footnotesize $r_R=27.6$};
  \draw[->] (-24,0) -- (24,0) node[right] {\footnotesize $x$ [m]};
  \draw[->] (-24,0) -- (-24,22) node[above] {\footnotesize $y$ [m]};
\end{tikzpicture}
\caption{Geometry used in Q2.3 (not to scale). The blue lines show the two
paths to the position $x = -10$\,m, which differ by almost exactly half a
wavelength.}
\label{fig:geometry}
\end{figure}

\paragraph{Assumptions.} Both the stage and the FOH tent are acoustically
transparent at 20\,Hz, and each subwoofer radiates as a monopole, so the
sound pressure falls as $1/r$ and its phase as $e^{-jkr}$. The subwoofers
and the listeners are all close to the ground, which is far less than a
wavelength at 20\,Hz. The ground reflection therefore adds in phase with the
direct sound everywhere and raises all levels by the same amount. Since the
subwoofers are calibrated by the level they produce at FOH, this common
factor is already included in the 80\,dB, and the problem can be treated as
free-field radiation from two monopoles.

\paragraph{Wavenumber and source strength.} At 20\,Hz
\begin{equation}
  k = \frac{2\pi f}{c} = \frac{2\pi\cdot 20}{343} = \SI{0.366}{rad/m},
  \qquad \lambda = \frac{c}{f} = \SI{17.15}{m}.
\end{equation}
The complex RMS pressure from a monopole is $p(r) = A\,e^{-jkr}/r$. Both
subwoofers are $r_\text{FOH} = \sqrt{10^2+13^2} = \SI{16.40}{m}$ from FOH and
each gives 80\,dB there, so they have the same strength:
\begin{equation}
  p_\text{FOH} = p_\text{ref}\,10^{80/20} = \SI{0.2}{Pa},\qquad
  |A| = p_\text{FOH}\,r_\text{FOH} = \SI{3.28}{Pa.m}.
\end{equation}
Both arrive at FOH with the same path length, i.e.\ in phase, so together
they would give 86\,dB there.

\paragraph{Coherent sum.} Both subwoofers are driven by the same 20\,Hz
signal, so their pressures are coherent and must be added as complex
numbers, not as energies:
\begin{equation}
  p = A\left(\frac{e^{-jkr_L}}{r_L} + \frac{e^{-jkr_R}}{r_R}\right),
  \qquad
  L_p = 20\log_{10}\frac{|p|}{p_\text{ref}},
\end{equation}
with $r_{L,R} = \sqrt{(x \pm 10)^2 + 19^2}$. The level from each subwoofer on
its own is $L = 80 + 20\log_{10}(r_\text{FOH}/r)$. The result depends on the
phase difference $k\,\Delta r$, where $\Delta r = |r_L - r_R|$ is the path
difference. Table~\ref{tab:q23} gives the numbers.

\begin{table}[htbp]
\centering
\caption{SPL at the five audience positions (dB re \SI{20}{\micro Pa}).
$L_L$, $L_R$: each subwoofer alone. Last column: energy sum, i.e.\ the
level without interference, for comparison.}
\label{tab:q23}
\begin{tabular}{rcccccc>{\bfseries}cc}
\toprule
$x$ [m] & $r_L$ [m] & $r_R$ [m] & $L_L$ & $L_R$ & $\Delta r$ [m] & $k\Delta r$ [$^\circ$] & $L_p$ & energy sum\\
\midrule
$-20$ & 21.47 & 35.51 & 77.7 & 73.3 & 14.04 & 295 & 80.4 & 79.0\\
$-10$ & 19.00 & 27.59 & 78.7 & 75.5 &  8.59 & 180 & 68.6 & 80.4\\
$0$   & 21.47 & 21.47 & 77.7 & 77.7 &  0.00 &   0 & 83.7 & 80.7\\
$10$  & 27.59 & 19.00 & 75.5 & 78.7 &  8.59 & 180 & 68.6 & 80.4\\
$20$  & 35.51 & 21.47 & 73.3 & 77.7 & 14.04 & 295 & 80.4 & 79.0\\
\bottomrule
\end{tabular}
\end{table}

\answer{SPL at $x = -20, -10, 0, 10, 20$\,m: \quad
\textbf{80.4, 68.6, 83.7, 68.6, 80.4\,dB} re \SI{20}{\micro Pa}.}

\paragraph{Discussion.} On the centre line the two paths are equal, so the
subwoofers add in phase: $77.7 + 6.0 = \SI{83.7}{dB}$. At $x = \pm\SI{10}{m}$
the path difference is \SI{8.59}{m}, almost exactly $\lambda/2$, so the two
pressures are in antiphase. They do not cancel completely only because the
nearer subwoofer is $3.2$\,dB louder than the farther one, which leaves a dip
of \SI{68.6}{dB}, about 15\,dB below the centre. At $x = \pm\SI{20}{m}$ the
phase difference is $295^\circ$ ($-65^\circ$), close enough to a full cycle
that the waves partly reinforce each other again. Figure~\ref{fig:line}
shows the full level variation along the audience line, and
Figure~\ref{fig:map} shows the interference pattern in front of the stage.
It is the classic left--right subwoofer ``power alley'' on the centre line,
with deep notches either side of it.

\begin{figure}[htbp]
\centering
\includegraphics[width=0.85\linewidth]{spl_audience_line.pdf}
\caption{SPL along the audience line $y = \SI{19}{m}$ at 20\,Hz. The solid
line is the coherent sum, the dashed line is the energy sum without
interference, and the dots are the five positions asked for.}
\label{fig:line}
\end{figure}

\begin{figure}[htbp]
\centering
\includegraphics[width=0.85\linewidth]{spl_map.pdf}
\caption{SPL map at 20\,Hz in front of the stage. Squares: subwoofers;
triangle: FOH; circles: audience positions.}
\label{fig:map}
\end{figure}

\clearpage
\appendix
\section*{Appendix --- MATLAB scripts}

\subsection*{problem1.m}
\lstinputlisting{../matlab/problem1.m}

\subsection*{problem2.m}
\lstinputlisting{../matlab/problem2.m}

\end{document}
